\documentclass[a4paper]{article}
% Español (México) con XeLaTeX: fontspec para fuentes Unicode, polyglossia
% para la división silábica y los nombres del documento en la variante mexicana.
\usepackage{fontspec}
\usepackage{polyglossia}
\setdefaultlanguage[variant=mexican]{spanish}
% Los nombres chinos van como «pinyin (original)»: xeCJK da a los caracteres
% Han su propia fuente; sin ella XeLaTeX los omite con un aviso.
% FandolSong no cubre algunos caracteres de nombres (U+73FA); WenQuanYi Zen Hei sí.
\usepackage[AutoFallBack=true]{xeCJK}
\setCJKmainfont{FandolSong-Regular.otf}
\setCJKfallbackfamilyfont{\CJKrmdefault}{WenQuanYi Zen Hei}
% polyglossia no traduce los nombres de listings.
\AtBeginDocument{\providecommand{\lstlistingname}{}\renewcommand{\lstlistingname}{Listado}%
  \providecommand{\lstlistlistingname}{}\renewcommand{\lstlistlistingname}{Índice de listados}}
\usepackage{amsmath, amssymb}
\usepackage{graphicx}
\usepackage[margin=2.5cm]{geometry}
\usepackage[most]{tcolorbox}
\usepackage{etoolbox}
\usepackage{listings}
\usepackage{hyperref}
\usepackage{booktabs}
\usepackage{subcaption}
\usepackage{float}

\newtcolorbox{knowledgebox}[1]{enhanced, colback=blue!5!white, colframe=blue!75!black, colbacktitle=blue!75!black, coltitle=white, fonttitle=\bfseries, title={#1}, attach boxed title to top left={yshift=-2mm, xshift=2mm}, boxrule=1pt, sharp corners}
\newtcolorbox{importantbox}[1]{enhanced, colback=yellow!10!white, colframe=yellow!80!black, colbacktitle=yellow!80!black, coltitle=black, fonttitle=\bfseries, title={#1}, sharp corners}
\newtcolorbox{warningbox}[1]{enhanced, colback=red!5!white, colframe=red!75!black, colbacktitle=red!75!black, coltitle=white, fonttitle=\bfseries, title={#1}, sharp corners}

\lstset{language=Python, basicstyle=\ttfamily\small, keywordstyle=\color{blue}, stringstyle=\color{red!60!black}, commentstyle=\color{green!60!black}, breaklines=true, frame=single, numbers=left, numberstyle=\tiny\color{gray}, captionpos=b, extendedchars=true}

\newcommand{\notetitle}{{[}LLM Agents F24{]} Project GR00T: A Blueprint for Generalist Robotics — Jim Fan}
\newcommand{\noteauthors}{Elaboradas a partir de materiales públicos del curso}
\newcommand{\notedate}{4 de noviembre de 2024}
\newcommand{\videochannel}{Berkeley RDI}
\newcommand{\videopublishdate}{4 de noviembre de 2024}
\newcommand{\videoduration}{aproximadamente 70 minutos}
\newcommand{\videourl}{https://www.youtube.com/live/Qhxr0uVT2zs}
\newcommand{\videocoverpath}{cover.jpg}
\newcommand{\repourl}{https://github.com/hqhq1025/ai-course-notes}


% Footer with repo info
\usepackage{fancyhdr}
\pagestyle{fancy}
\fancyhf{}
\fancyfoot[C]{\thepage}
\fancyfoot[R]{\footnotesize\color{gray}\href{\repourl}{AI Course Notes}}
\renewcommand{\headrulewidth}{0pt}
\renewcommand{\footrulewidth}{0pt}

\begin{document}
\begin{titlepage}\centering
{\Large Notas del curso\par}\vspace{1.2cm}
{\huge\bfseries \notetitle\par}\vspace{0.8cm}
{\large \noteauthors\par}\vspace{0.3cm}
{\large \notedate\par}\vspace{1.2cm}
\ifdefempty{\videocoverpath}{{\small Favor de llenar la ruta de la imagen de portada.\par}}{\includegraphics[width=0.82\textwidth,height=0.45\textheight,keepaspectratio]{\videocoverpath}\par}
\vfill
\begin{tcolorbox}[width=0.9\textwidth, colback=black!2!white, colframe=black!60, sharp corners]
\textbf{Autor o canal del video}: \videochannel\par
\textbf{Fecha de publicación}: \videopublishdate\par
\textbf{Duración del video}: \videoduration\par
\ifdefempty{\videourl}{}{\textbf{Enlace del video}: \href{\videourl}{\nolinkurl{\videourl}}\par}\par
\textbf{Página del proyecto}: \href{\repourl}{\nolinkurl{\repourl}}\par
\end{tcolorbox}
\end{titlepage}

\tableofcontents\newpage

\section{Introducción: la historia de los dos gatitos}

Jim Fan es científico investigador sénior de NVIDIA y dirige Project GR00T (modelo fundacional de robótica general). Abrió la clase con el experimento clásico de Held y Hein de 1963: dos gatitos recién nacidos en un dispositivo de carrusel, donde el gatito activo podía moverse libremente y observar el mundo, mientras el gatito pasivo solo podía mirar de manera pasiva; como resultado, solo el gatito activo desarrolló un sistema visual sano.

\begin{importantbox}{Metáfora central: la interacción incorporada es la base de la inteligencia}
La observación pasiva no basta para producir una comprensión verdadera: la inteligencia requiere \textbf{interacción activa} con el mundo físico. Esta es precisamente la diferencia fundamental entre los LLM (el gatito pasivo, que solo lee texto de internet) y los Agentes incorporados (el gatito activo, que interactúa con el mundo físico).
\end{importantbox}

\subsection{Resumen de la sección}

La AGI verdadera no puede lograrse solo leyendo texto; requiere inteligencia incorporada (embodied intelligence): el Agente debe poder percibir, actuar y aprender a partir de la retroalimentación del mundo físico.

\section{Foundation Agent: de lenguaje a acción}

\subsection{Aplicación de los modelos base en los Agentes}

Jim Fan propone el concepto de \textbf{Foundation Agent}: utilizar modelos base preentrenados a gran escala para construir Agentes de propósito general:

\begin{itemize}
    \item \textbf{Voyager}: un Agente de exploración de mundo abierto en Minecraft que usa GPT-4 como núcleo de razonamiento, descubre habilidades automáticamente, escribe código y lo almacena en una biblioteca de habilidades, logrando aprendizaje continuo
    \item \textbf{Eureka}: usa un LLM para generar y optimizar automáticamente la reward function, entrenando manipulación diestra con manos en simulación física
    \item \textbf{MineDojo}: una base de conocimiento de Minecraft a gran escala y un banco de pruebas de mundo abierto
\end{itemize}

\begin{knowledgebox}{La biblioteca de habilidades de Voyager}
Cada habilidad que Voyager aprende (como fabricar una espada o construir una casa) se almacena en forma de código ejecutable. Las tareas posteriores pueden invocar directamente las habilidades ya existentes, sin necesidad de aprender desde cero; este es un excelente ejemplo de memoria a largo plazo y aprendizaje continuo en los Agentes.
\end{knowledgebox}

\subsection{Resumen de la sección}

Mediante las capacidades de razonamiento y generación de código de los LLM, Foundation Agent logra exploración de mundo abierto y aprendizaje continuo de habilidades en mundos virtuales.

\section{Project GR00T: modelo fundacional general para robots}

\subsection{Visión y arquitectura}

El objetivo de GR00T es construir un modelo fundacional general capaz de operar diversos cuerpos robóticos:

\begin{itemize}
    \item Recibe instrucciones en lenguaje natural y entradas visuales
    \item Genera secuencias de acciones robóticas
    \item Generaliza entre distintas morfologías de robots: el mismo modelo controla robots diferentes
\end{itemize}

\subsection{Retos técnicos centrales}

\begin{warningbox}{Retos particulares de los modelos fundacionales para robots}
\begin{itemize}
    \item \textbf{Escasez de datos}: a diferencia del lenguaje, que cuenta con datos a escala de internet, los datos de operación robótica son extremadamente escasos
    \item \textbf{sim-to-real gap}: las políticas entrenadas en simulación suelen fallar en el mundo real
    \item \textbf{Seguridad}: los robots actúan directamente sobre el mundo físico, y un error puede causar daño físico
    \item \textbf{Tiempo real}: la interacción física exige una latencia de decisión del orden de milisegundos
\end{itemize}
\end{warningbox}

\subsection{Estrategias de generación de datos}

\begin{itemize}
    \item Utiliza simulación física a gran escala (Isaac Sim) para generar datos de entrenamiento
    \item Genera automáticamente tareas diversas y reward function mediante LLM (método Eureka)
    \item Aprende operaciones humanas a partir de videos de internet (aprendizaje por observación)
    \item Human demonstration + teleoperation para recopilar datos reales de operación
\end{itemize}

\subsection{Resumen de la sección}

GR00T representa la ambición de extender el paradigma de los modelos fundacionales del lenguaje al mundo físico; los retos centrales están en los datos, la transferencia sim-to-real y la seguridad.
\section{Arquitectura del sistema y ruta de datos}

\subsection{Niveles de percepción multimodal}

La arquitectura de GR00T está compuesta por tres niveles de percepción visual/lingüística/de fuerza:

\begin{itemize}
    \item \textbf{Codificador visual}: procesa RGB-D y escaneos de lidar, y produce un embedding de percepción espacial
    \item \textbf{Programador de lenguaje}: analiza las instrucciones para obtener la intención de la acción, la secuencia de llamadas a herramientas y alinea el reward
    \item \textbf{Programador de acciones}: selecciona torque/force/action según las restricciones dinámicas, en un ciclo cerrado en tiempo real
\end{itemize}

\begin{knowledgebox}{Silogismo percepción-lenguaje-acción}
Estos tres niveles forman un ciclo cerrado de percepción-intención-ejecución: el módulo visual le dice al módulo de lenguaje «qué hay en la escena», el módulo de lenguaje le dice al módulo de acción «qué quiero hacer», y el módulo de acción retroalimenta el resultado y actualiza la política según el reward.
\end{knowledgebox}

\subsection{Flujo de datos y ruta de entrenamiento}

La ruta de entrenamiento típica es la siguiente:

\begin{enumerate}
    \item Recopilar datos de teleoperation/telemetry (operación humana + sensor log)
    \item Generar automáticamente el prompt de la tarea con un LLM (Eureka pipeline)
    \item Expandir la tarea en Isaac Sim y generar un synthetic replay
    \item Aplicar offline RL sobre el replay (off-policy data + behavior cloning)
    \item Aplicar sim-to-real gap alignment (domain randomization + system ID)
\end{enumerate}

\begin{importantbox}{Replay Buffer hygiene}
Mantener diversidad en el replay buffer: trajectories exitosas recientes, unas cuantas muestras fallidas, expert demonstration, LLM-generated exploration. Combinar el replay secuencial con el prioritized replay evita que el modelo sobreajuste las muestras comunes.
\end{importantbox}

\subsection{Detalles de Sim-to-Real Alignment}

La estrategia de puente de GR00T adopta una triple garantía:

\begin{enumerate}
    \item \textbf{Domain randomization}: perturba la iluminación, la textura y el coeficiente de fricción para que el modelo se adapte a escenarios físicos más amplios
    \item \textbf{System identification}: utiliza una pequeña cantidad de real-world trajectory para ajustar los parámetros de la simulación y asegurar que la dinámica corresponda
    \item \textbf{Validation loop}: ejecuta tareas simplificadas en un robot real y registra el número de casos fuera de distribución en el DevOps dashboard
\end{enumerate}

\begin{importantbox}{Sim-to-Real leash}
Cada deployment viene acompañado de un ``sim-to-real leash'': mientras el Validation loop no se apruebe, el modelo se mantiene en modo sandbox, y solo se libera cuando se detecta un tolerable drift level.
\end{importantbox}

\subsection{Resumen de la sección}

Sim-to-Real alignment requiere que domain randomization, system identification y validation loop formen un ciclo cerrado entre los tres; el leash y el guard rail garantizan una transición progresiva entre la simulación y la realidad.

\section{Casos y operación: implementación de GR00T fuera de Berkeley}

\subsection{Resumen de casos de la industria}

Jim Fan analizó dos casos reales:

\begin{itemize}
    \item \textbf{Laboratorio de obleas de semiconductores}: usa el modelo GR00T junto con un brazo robótico de sujeción autónomo para completar el intercambio de wafer. Se apoya en built-in alignment rules para evitar colisiones.
    \item \textbf{Centro de clasificación logística}: el objetivo es recoger objetos de formas diversas en un entorno abarrotado, generando planes de ejecución mediante instrucciones de LLM y detectando errores mediante per-instance vision.
\end{itemize}

\begin{warningbox}{La incomodidad real al implementar}
En los escenarios reales, la falla más común no es un error de inferencia del modelo, sino sensor calibration drift, el backlash mecánico y las pequeñas desviaciones que provoca el cambio de herramienta. Incluso el modelo más potente necesita un fallback confiable.
\end{warningbox}

\subsection{Métricas de operación y monitoreo}

Para garantizar la implementación, la cadena de observación se divide en tres capas:

\begin{itemize}
    \item \textbf{Trace}: cada acción/instrucción incluye timestamp, la secuencia de tokens y el tool call log
    \item \textbf{Metrics}: tasa de éxito, cycle time, safety intervention frequency
    \item \textbf{Alerts}: los thresholds preestablecidos (colisión, tiempo de espera agotado de la tarea, planner oscillation) activan la luz roja al dispararse
\end{itemize}

\begin{knowledgebox}{Trace-driven postmortem}
Guardar el trace buffer permite reproducir rápidamente lo ocurrido después de un incident: cada token, cada llamada a herramienta y cada transición de estado puede rastrearse, lo que ayuda a los ingenieros a localizar «qué instrucción del paso 37 provocó la desviación».
\end{knowledgebox}

\subsection{Flujo operativo y Playbook}

Cada proyecto de manipulación sigue un Operational Playbook de cinco pasos:

\begin{enumerate}
    \item \textbf{Plan}: definir la tarea, configurar el documento CTA y generar failure hypotheses
    \item \textbf{Deploy}: escribir la live config, calibrar los sensors e insertar el override switch
    \item \textbf{Observe}: monitorear drop rates, collisions y latency mediante metrics + open telemetry
    \item \textbf{Adjust}: si se dispara una alert, entrar en modo «pause + analyze», reproducir el trace y ajustar la instruction
    \item \textbf{Document}: escribir un postmortem por cada incident y actualizar el Checklist
\end{enumerate}

\begin{warningbox}{Pause + Analyze es el freno más seguro}
Cuando el desempeño se degrada, no ajuste el large language predicate a ciegas; basta con seguir el checklist para poner en pausa el robot, reproducir el trace y revisar el lockstep metric para evitar un cascade failure.
\end{warningbox}

\subsection{Resumen de la sección}

La estabilidad operativa requiere el respaldo de Case study, monitoreo/trace/backstop, y debe combinarse con el playbook para poder enfrentar los incident de la implementación real.

\newpage

\section{Aspectos destacados de las diapositivas}

\subsection{Extractos visuales}

Jim Fan enumeró en su guion una serie de diapositivas clave; a continuación se seleccionan páginas representativas para su lectura:

\begin{figure}[p]
\centering
\includegraphics[width=0.95\textwidth,page=2]{slides.pdf}
\caption{Marco semántico y biblioteca de habilidades de Foundation Agent.}
\end{figure}

\begin{figure}[p]
\centering
\includegraphics[width=0.95\textwidth,page=3]{slides.pdf}
\caption{Pipeline de percepción multimodal de GR00T (Vision + Language + Actuation).}
\end{figure}

\begin{figure}[p]
\centering
\includegraphics[width=0.95\textwidth,page=5]{slides.pdf}
\caption{Esquema de domain randomization y guard rail para el Sim-to-Real alignment.}
\end{figure}

\begin{figure}[p]
\centering
\includegraphics[width=0.95\textwidth,page=7]{slides.pdf}
\caption{Operational Playbook y flujo de reproducción de incident trace.}
\end{figure}

\begin{figure}[p]
\centering
\includegraphics[width=0.95\textwidth,page=9]{slides.pdf}
\caption{Ejemplo de Governance checklist y usage policy.}
\end{figure}

\begin{figure}[p]
\centering
\includegraphics[width=0.95\textwidth,page=11]{slides.pdf}
\caption{Multi-robot stack de Project GR00T y su collab plan.}
\end{figure}

\begin{figure}[p]
\centering
\includegraphics[width=0.95\textwidth,page=12]{slides.pdf}
\caption{Ilustración del pipeline que combina LLM + physics solver.}
\end{figure}

\begin{figure}[p]
\centering
\includegraphics[width=0.95\textwidth,page=14]{slides.pdf}
\caption{Replay buffer routing y prioritize schema.}
\end{figure}

\begin{figure}[p]
\centering
\includegraphics[width=0.95\textwidth,page=17]{slides.pdf}
\caption{Trace-driven observability board, combinado con metrics dashboard.}
\end{figure}

\begin{figure}[p]
\centering
\includegraphics[width=0.95\textwidth,page=19]{slides.pdf}
\caption{Detalles del flujo de trabajo Pause + analyze y safety alert icons.}
\end{figure}

\begin{figure}[p]
\centering
\includegraphics[width=0.95\textwidth,page=23]{slides.pdf}
\caption{Future research roadmap: LLM physics, self-supervised, multi-agent.}
\end{figure}

\begin{figure}[p]
\centering
\includegraphics[width=0.95\textwidth,page=25]{slides.pdf}
\caption{Incident response sequence: detect, pause, analyze, recover.}
\end{figure}

\begin{figure}[p]
\centering
\includegraphics[width=0.95\textwidth,page=27]{slides.pdf}
\caption{Reward decomposition chart showing language + physics components.}
\end{figure}

\begin{figure}[p]
\centering
\includegraphics[width=0.95\textwidth,page=30]{slides.pdf}
\caption{Logging taxonomy: sensors, policies, overrides.}
\end{figure}

\begin{figure}[p]
\centering
\includegraphics[width=0.95\textwidth,page=32]{slides.pdf}
\caption{Benchmark ladder comparing dexterity, reliability, generalization.}
\end{figure}

\begin{figure}[p]
\centering
\includegraphics[width=0.95\textwidth,page=34]{slides.pdf}
\caption{Toolchain health dashboard with safety interlocks.}
\end{figure}

\begin{figure}[p]
\centering
\includegraphics[width=0.95\textwidth,page=35]{slides.pdf}
\caption{Data lineage schematic showing robot teleop -> replay -> offline RL.}
\end{figure}

\begin{figure}[p]
\centering
\includegraphics[width=0.95\textwidth,page=36]{slides.pdf}
\caption{Governance maturity curve with audit trail metrics.}
\end{figure}

\begin{figure}[p]
\centering
\includegraphics[width=0.95\textwidth,page=37]{slides.pdf}
\caption{Monitoring dashboard outlining trace, metrics, alerts.}
\end{figure}

\begin{figure}[p]
\centering
\includegraphics[width=0.95\textwidth,page=38]{slides.pdf}
\caption{Incident review template with pause/analyze metrics.}
\end{figure}

\begin{figure}[p]
\centering
\includegraphics[width=0.95\textwidth,page=39]{slides.pdf}
\caption{Lifecycle audit trail showing dataset, config, release pairings.}
\end{figure}

\subsection{Resumen de la sección}

Las diapositivas visualizan la arquitectura de GR00T, la estrategia de simulación, el playbook operativo y el checklist de gobernanza; son material importante para concretar el contenido del texto.

\newpage

\section{Gobernanza y trabajo futuro}

\subsection{Marco de gobernanza}

GR00T no se limita a la investigación, también busca alinearse con la regulación:

\begin{itemize}
    \item \textbf{Safety board}: antes de cada release debe haber una evaluación humana de los reward escalation scenarios
    \item \textbf{Usage policy}: define qué tareas pueden ejecutarse de forma autonomous y cuáles deben ser human-in-loop
    \item \textbf{Audit trail}: registra de forma completa el dataset, los hyperparam y el model ID de cada deployment
\end{itemize}

\begin{importantbox}{Lista de verificación de governance de nivel industrial}
1. ¿Existe un failure mode radical (por ejemplo, tomar un objeto punzocortante)? 2. ¿Se proporciona un override switch? 3. ¿Se registran la instrucción y la response? 4. ¿Existe una alerta de non-nominal behavior?
\end{importantbox}

\subsection{Líneas futuras de investigación}

Jim Fan plantea tres caminos:

\begin{enumerate}
    \item \textbf{LLM + physics solvers}: incorporar differentiable physics a la planeación
    \item \textbf{Self-supervised manipulation}: sintetizar demonstrations de forma automática mediante observation consistency
    \item \textbf{Multi-agent robot swarms}: hacer que múltiples versiones de GR00T colaboren para completar procesos industriales
\end{enumerate}

\subsection{Resumen de la sección}

La gobernanza, la usage policy, el audit trail y las líneas futuras de investigación forman en conjunto la base de seguridad para la implementación industrial.

\section{El futuro de la IA corpórea}

\begin{importantbox}{El cero a un billón de dólares}
En palabras de Jensen Huang: la robótica de propósito general es ``la industria de un billón de dólares del futuro''—porque implica la automatización de todo el mercado laboral. Pero por ahora es una industria de ``cero billones'', porque todavía no existe un producto realmente utilizable. Precisamente porque se trata de pasar de cero a un billón, este es el mejor momento para invertir en investigación.
\end{importantbox}

Jim Fan está convencido de que la AGI no puede existir sin inteligencia corpórea—algún día construiremos, con materiales basados en silicio, ``gatitos activos''.

\subsection{Resumen de la sección}

La IA corpórea es una parte indispensable de la AGI. La tecnología actual está migrando del éxito en mundos simulados hacia el mundo físico.

\section{Síntesis y ampliación}

\subsection{Puntos clave}

\begin{enumerate}
    \item La interacción activa es la base de la inteligencia: un LLM de solo texto es un ``gatito pasivo''
    \item Foundation Agent (Voyager, Eureka) muestra la exploración de mundo abierto impulsada por LLM y el diseño automático de reward
    \item Project GR00T busca construir un modelo base general entre distintas morfologías de robots
    \item La escasez de datos y el sim-to-real gap son los retos centrales
    \item La IA con cuerpo es un componente necesario de la AGI
\end{enumerate}

\subsection{Summary table}

\begin{table}[H]
\centering
\begin{tabular}{p{0.24\textwidth} p{0.2\textwidth} p{0.28\textwidth} p{0.2\textwidth}}
\toprule
Tema & Promesa & Reto & Producto clave \\
\midrule
Interacción activa con cuerpo & Ciclo cerrado de percepción-lenguaje-acción a la medida del robot & Escasez de datos + sim-to-real gap & Arquitectura de Voyager/Eureka, baseline de GR00T \\
Datos y entrenamiento & LLM-generated tasks + replay hygiene & Mantener la diversidad, distribución natural & Isaac Sim pipelines, generación de reward con LLM \\
Operación en producción & Trace + metrics + alerts & Cambio de herramientas, calibration drift & Incident playback, checklists \\
Gobernanza & Safety board + audit trail & Autonomous failure modes & Usage policy + override switch \\
\bottomrule
\end{tabular}
\caption{Principales dimensiones de aplicación práctica de la Lección 04}
\end{table}

\subsection{Lecturas adicionales}

\begin{itemize}
    \item Wang et al., ``Voyager: An Open-Ended Embodied Agent with Large Language Models,'' 2023.
    \item Ma et al., ``Eureka: Human-Level Reward Design via Coding Large Language Models,'' ICLR 2024.
    \item Fan et al., ``MineDojo: Building Open-Ended Embodied Agents with Internet-Scale Knowledge,'' NeurIPS 2022.
\end{itemize}


\end{document}
